\documentclass[12pt,a4paper]{article}

% ===========================================================================
% prepare2.tex  —  DEFENSE Q&A CHEAT-SHEET (private prep, not part of thesis)
% ---------------------------------------------------------------------------
% Mỗi mục gồm:  Q = câu hội đồng có thể hỏi
%               A = câu trả lời mẫu: NGẮN, CHẮC, ĐÓNG CHỦ ĐỀ.
% Nguyên tắc để KHÔNG bị hỏi sâu hơn:
%   1. Trả lời 2–3 câu, nêu ngay "lý do thiết kế" rồi dừng.
%   2. Chủ động nhận một giới hạn nhỏ ("đây là lựa chọn có chủ đích / là future
%      work") — reviewer thấy mình đã lường trước thì không đào tiếp.
%   3. Không hứa những gì code không làm; không bịa số.
% ===========================================================================

% --- Vietnamese-safe setup: works under pdfLaTeX AND XeLaTeX/LuaLaTeX --------
% (pdfLaTeX needs T5 fontenc for Vietnamese glyphs; XeLaTeX/LuaLaTeX use
%  fontspec + a Unicode font. Overleaf: chọn engine nào cũng biên dịch được.)
\usepackage{iftex}
\ifPDFTeX
  \usepackage[utf8]{inputenc}
  \usepackage[T5]{fontenc}
\else
  \usepackage{fontspec}
\fi
\usepackage[vietnamese,english]{babel}
\usepackage[margin=1in]{geometry}
\usepackage{xcolor}
\usepackage{enumitem}
\usepackage{hyperref}
\hypersetup{colorlinks=true, linkcolor=black, urlcolor=blue}

% --- helper: một cặp Hỏi / Đáp -------------------------------------------
\newcommand{\QA}[2]{%
  \par\medskip\noindent\textbf{\textcolor{blue!60!black}{Q.}} #1\par
  \smallskip\noindent\textbf{\textcolor{green!45!black}{A.}} #2\par}

\newcommand{\tip}[1]{\par\smallskip\noindent%
  \textit{\textcolor{gray}{[mẹo: #1]}}\par}

\title{\bfseries Defense Q\&A --- Sổ tay chuẩn bị bảo vệ\\
       \large 3D House Defense Game}
\author{Nguyen Thuy Ngoc --- 23BI14344}
\date{Hanoi, June 2026}

\begin{document}
\maketitle

\noindent\textit{Tài liệu cá nhân để luyện trả lời. Mục tiêu: mỗi câu trả lời
gọn 2--3 ý, nêu \textbf{lý do thiết kế} rồi dừng lại. Nếu bị dồn, quy về một
trong ba câu chốt an toàn:}
\begin{itemize}[nosep]
  \item \textit{``Đây là lựa chọn có chủ đích để giữ phạm vi single-player.''}
  \item \textit{``Cái này em có ghi trong phần Limitations / Future Work.''}
  \item \textit{``Luật đó được quyết trong reducer thuần nên em test được, mời
        thầy xem \texttt{reducer.test.ts}.''}
\end{itemize}

\tableofcontents
\bigskip

% ===========================================================================
\section{Kiến trúc \& ``single source of truth''}
% ===========================================================================

\QA{Thầy thấy luận văn nói ``một state duy nhất'', nhưng spawn scheduler, collision
và renderer đều giữ state riêng. Có mâu thuẫn không?}
{Em phân biệt \emph{simulation state} và \emph{scheduling/derived state}. Toàn bộ
\emph{luật chơi} --- enemy, đạn, máu, wave, thắng/thua --- nằm trong một
\texttt{GameState} và chỉ đổi qua reducer. Scheduler chỉ giữ ``đã spawn tới con
thứ mấy'', renderer chỉ giữ map id$\to$mesh: đó là chi tiết thực thi, không phải
luật chơi, nên cố ý để ngoài để reducer thuần và test được. Nếu cần, các con số
roster vẫn nằm ở \texttt{waves.ts} để reducer tự phát hiện wave xong.}
\tip{nhấn ``luật chơi thì trong reducer'', ``lịch trình thì ngoài'' --- rồi dừng.}

\QA{Vì sao trong GameState vẫn có \texttt{player.hp = 100} mà không dùng?}
{Ban đầu em thiết kế thua theo máu nhà, sau đổi sang thua theo \emph{vị trí}
(enemy chạm vạch nhà là thua ngay) vì nó rõ ràng hơn cho người chơi. Trường
\texttt{hp} là tàn dư của thiết kế cũ; nó không ảnh hưởng luật hiện tại. Em sẽ dọn
nó ở bản cuối.}
\tip{nhận thẳng là tàn dư + sẽ dọn = đóng chủ đề, đừng cố bào chữa.}

\QA{Chống double-hit (một vụ nổ chỉ gây damage một lần) làm ngoài reducer bằng
một tập ``đã impact''. Sao không đưa vào reducer?}
{Vụ nổ là một \emph{sự kiện tức thời}, không phải trạng thái lâu dài, nên em xử lý
ở game loop và chỉ bắn ra các action \texttt{DAMAGE\_ENEMY} --- reducer vẫn là nơi
duy nhất trừ máu. Set ``đã impact'' chỉ để loop không phát lại cùng một vụ nổ ở
frame sau; nó không quyết định luật, chỉ khử lặp.}

% ===========================================================================
\section{Luật chơi \& cân bằng số}
% ===========================================================================

\QA{Big Shot ghi ``every 3'' nhưng đọc code thì phát thứ 4 mới là Big Shot?}
{Đúng, ý nghĩa là ``sau mỗi 3 phát thường thì có 1 phát mạnh''. Bộ đếm chạy
0,1,2 cho ba phát thường rồi phát kế là Big Shot và reset. Cách này để người chơi
luôn được thưởng đều đặn sau một nhịp bắn ổn định.}

\QA{Advanced bắn 2 đạn, blast bán kính 6 mà chỉ cách nhau 1.5 --- một enemy có thể
ăn damage từ cả hai. Cố ý hay lỗi?}
{Cố ý. Đây là phần thưởng của tier cao nhất: nếu người chơi ngắm khéo vào một cụm,
hai vụ nổ chồng lên cho damage cộng dồn --- đó chính là lý do Advanced đáng để
mở khóa. Damage vẫn có trần vì mỗi enemy có máu giới hạn.}
\tip{gọi nó là ``reward for good aim'', đừng gọi là ``bug''.}

\QA{Người chơi có thể chọn lại vũ khí yếu hơn đã mở. Có phải lỗi?}
{Không. Em cho tự do đổi giữa mọi vũ khí đã mở vì mỗi tier có đánh đổi riêng ---
Basic reload nhanh hơn, đạn khác tốc độ. Cho người chơi chọn theo tình huống làm
phần thưởng ``được chọn'' chứ không bị ép, đúng như thay đổi so với SRS gốc mà em
nêu trong luận văn.}

\QA{Thời gian băng đồng của enemy là bao nhiêu --- code có chỗ ghi 8s/15s?}
{Con số chuẩn là 25/30/35 giây (easy/medium/hard), đúng như bảng trong luận văn.
Dòng 8s/15s là comment cũ từ lúc prototype, em quên cập nhật; giá trị thực trong
bảng \texttt{TRAVEL\_TIME} là 25/30/35. Em đã tune chậm lại sau khi chơi thử,
điều này em có ghi ở chương Kết quả.}
\tip{chốt ngay ``25/30/35 là chuẩn'' để hội đồng không lấy số sai.}

\QA{``25 giây băng đồng'' là theo trục thẳng hay theo đường đi lượn sóng?}
{Theo trục tiến về nhà (trục x) --- tốc độ theo trục đó là hằng số nên thời gian
tới nhà luôn đúng 25/30/35s bất kể đường lượn. Đường lượn sin chỉ để enemy tản ra
theo chiều rộng cho đỡ dồn một hàng, không làm chậm tiến độ về nhà.}

\QA{Vì sao thua theo vị trí, không theo máu nhà?}
{Vì mục tiêu người chơi phải cực rõ: ``đừng để con nào lọt qua''. Thua ngay khi
một con chạm vạch nhà giữ được độ căng, và em kiểm tra được luật này bằng test
thuần trong reducer.}

% ===========================================================================
\section{Backend, dữ liệu \& leaderboard}
% ===========================================================================

\QA{Vì sao có tới ba lớp chống trùng submit (ref, map 30s, DB 10s)?}
{Ba lớp ở ba mức khác nhau: ref chặn React Strict Mode gọi effect hai lần trong
cùng một lần chơi; map trên client chặn double-click gửi lại; cửa sổ 10s ở DB là
lưới an toàn cuối cùng nếu request tới hai lần do mạng. Chúng bổ trợ nhau để đảm
bảo ``một lần chơi = một bản ghi''.}
\tip{nếu bị hỏi về false-positive (hai người trùng hệt trong 10s): ``xác suất cực
thấp với tên + thời gian mili-giây; em chấp nhận đánh đổi này để chống double-post,
và có thể siết bằng token mỗi ván ở future work''.}

\QA{API leaderboard mở, không auth --- làm sao chống điểm giả?}
{Trong phạm vi single-player, cục bộ, em ưu tiên tính đơn giản: server chỉ
\emph{validate hình dạng} dữ liệu (tên, thời gian $>0$, wave 1--3, won boolean).
Chống gian lận thật sự cần tài khoản và token phía server --- em xếp vào Future
Work cùng với việc thêm player accounts.}
\tip{đừng cố bảo vệ như thể đã chống gian lận; quy thẳng về Future Work.}

\QA{Bốn khóa sắp xếp nhưng với winner thì wave luôn 3, với loser thì bỏ time ---
vậy có thừa không?}
{Không thừa: bốn khóa xử lý \emph{mọi} trường hợp trong một câu \texttt{ORDER BY}.
Với nhóm winner, khóa quyết định là thời gian; với nhóm chưa thắng, là wave sâu
nhất rồi tới mới nhất. Viết đủ bốn khóa để một truy vấn duy nhất trả về đúng thứ
hạng cho cả hai nhóm mà client không phải sắp xếp lại.}

\QA{SQLite file sẽ không tồn tại trên host stateless --- vậy demo chạy ở đâu?}
{Em demo trên máy cục bộ (localhost), nơi file SQLite hoạt động đầy đủ. Em có nêu
rõ trong Limitations rằng để deploy công khai cần thay bằng DB có host; nhờ mọi
truy cập DB gói trong một file \texttt{db.ts}, việc đổi sang Postgres chỉ sửa
đúng file đó.}

% ===========================================================================
\section{Hiệu năng \& đo đạc}
% ===========================================================================

\QA{Sao chương Kết quả không có ảnh chụp màn hình game / leaderboard?}
{Em cố ý không đưa screenshot tĩnh, vì thứ quan trọng nhất của game này là
\emph{chuyển động} --- di chuyển real-time, ngắm bắn, phản hồi hình ảnh --- ảnh
tĩnh không thể hiện được. Thay vào đó em trình diễn \textbf{trực tiếp} ngay sau
phần thuyết trình: đi trọn một ván từ welcome, countdown, combat, Big Shot, phần
thưởng đổi vũ khí, tới màn thắng, và leaderboard cập nhật theo các ván vừa chơi.}
\tip{nói ngay ``em sẽ demo trực tiếp'' rồi mời hội đồng xem cuối buổi --- đóng chủ đề.}

\QA{Số FPS đo bằng gì, trên bao nhiêu máy?}
{Em đo bằng Chrome DevTools (Performance / FPS meter) trên một máy desktop, trong
lúc chơi thật ở wave đông nhất. Kết quả: đa số thời gian chạm trần 60\,Hz, chỉ tụt
ngắn xuống khoảng 31 khi combat nặng --- luôn $\ge$ mục tiêu 30 FPS. Em nêu rõ đây
là một cấu hình tham chiếu, không phải benchmark nhiều máy.}
\tip{chủ động nói ``một máy, một cấu hình tham chiếu'' để không bị hỏi ``sao ít mẫu''.}

\QA{Vì sao khẳng định dip FPS là do garbage collection?}
{Đó là giải thích hợp lý nhất: mô hình state bất biến tạo nhiều object ngắn hạn mỗi
frame, GC thu hồi chúng gây khựng ngắn. Em chưa profile allocation chi tiết nên nói
đây là nguyên nhân \emph{chính khả dĩ}; hướng tối ưu (tái dùng object, instanced
mesh) em để ở Future Work vì FPS chưa bao giờ tụt dưới mục tiêu.}
\tip{dùng chữ ``khả dĩ / likely'', đừng khẳng định tuyệt đối nguyên nhân.}

\QA{Delta-time bị chặn ở 1/15s có ý nghĩa gì?}
{Để một tab chạy nền quay lại không làm simulation nhảy cóc một bước lớn. Khi frame
quá dài, em coi như một frame hơi chậm thay vì tua nhanh. Trên máy rất yếu điều này
làm mô phỏng chậm lại đôi chút --- em chọn ưu tiên tính ổn định/đúng đắn hơn là
đúng thời gian tuyệt đối, hợp lý cho một game single-player.}

% ===========================================================================
\section{Lý thuyết / đồ họa}
% ===========================================================================

\QA{Vì sao damage giảm tuyến tính theo khoảng cách, không dùng inverse-square vật lý?}
{Vì đây là lựa chọn \emph{gameplay}, không phải mô phỏng vật lý. Tuyến tính dễ hiểu
và công bằng: tâm nổ full damage, ra tới bán kính thì bằng 0, người chơi ước lượng
được ngay nên khỏi cần tay. Inverse-square sẽ dồn hầu hết sát thương vào một điểm
rất nhỏ, khó chơi hơn.}

\QA{Raycasting cắt với mặt phẳng toán học $y=0$ thay vì mesh sàn --- khác gì?}
{Cắt với một mặt phẳng toán học cho giao điểm chính xác và rẻ (một phép giải đại
số), không phụ thuộc lưới tam giác hay độ mấp mô của sàn hiển thị. Nhờ vậy điểm
ngắm luôn nằm đúng dưới con trỏ, ổn định bất kể mô hình sàn thay đổi.}

\QA{Đạn giải ngược vận tốc đứng để chạm đích tại $t=T$; cung có bị dựng đứng khi
bắn xa với thời gian bay ngắn không?}
{Về lý thuyết vận tốc đứng tăng khi $T$ nhỏ, nhưng ba tier có $T$ = 1.3/1.0/0.7s
được chọn sao cho cung vẫn nhìn tự nhiên trong kích thước sân thực tế. Em kiểm tra
bằng mắt khi tune và thấy ổn ở mọi điểm hợp lệ trong yard.}

% ===========================================================================
\section{Câu hỏi ``analyst'' (nửa phân tích của luận văn)}
% ===========================================================================

\QA{SRS đầy đủ đâu? Sao code tham chiếu FR/BR/TR mà luận văn không in ra danh mục?}
{Em trình bày yêu cầu dưới dạng \emph{quyết định thiết kế có kèm tiêu chí chấp
nhận} trong Chương 3, và cô đọng danh mục đại diện (kèm tiêu chí + vị trí code) ở
Phụ lục A. Bộ định danh FR/BR/TR là hệ đánh số em dùng xuyên suốt để truy vết từ
luật tới dòng code. Nếu thầy cần, em có bản danh mục chi tiết hơn để trình.}
\tip{nếu chưa có file SRS rời: nói ``bản đầy đủ em quản lý riêng, phần đại diện đã ở
Phụ lục A'' --- đừng hứa đưa cái không có.}

\QA{Em học phương pháp đặc tả (viết SRS) từ đâu?}
{Từ kỳ thực tập: em quan sát cách hệ thống được phân tích và đặc tả trước khi làm,
rồi áp dụng vào đồ án độc lập này bằng cách viết yêu cầu có tiêu chí kiểm thử. Đây
là lần đầu em tự thực hành trọn vẹn nửa phân tích, nên bài học lớn nhất là ``một
yêu cầu chỉ hữu ích khi kiểm chứng được bằng cách chơi''.}

\QA{Ba thay đổi so với SRS gốc là gì và vì sao?}
{Một, cho người chơi \emph{chọn} có dùng vũ khí mới hay không (thưởng thấy xứng
đáng). Hai, bỏ điểm số trên HUD (wave + thời gian đã đủ, điểm gây rối). Ba, tăng
thời gian băng đồng chậm lại sau khi chơi thử vì nhịp gốc quá gấp. Cả ba đều là
quyết định lúc play-test và em ghi lại để thiết kế trung thực.}

% ===========================================================================
\section{Ba câu chốt an toàn (học thuộc)}
% ===========================================================================

\begin{enumerate}[leftmargin=*]
  \item \textbf{Khi bị hỏi ``sao không làm X''}: ``Phạm vi đồ án là single-player,
        desktop, chuột --- X (multiplayer / account / mobile) em cố ý để ngoài và
        nêu ở Future Work.''
  \item \textbf{Khi bị hỏi về một con số}: quy về nguồn duy nhất ---
        ``số đó nằm ở \texttt{constants.ts} / \texttt{weapons.ts} / \texttt{waves.ts},
        mỗi giá trị chỉ khai báo một chỗ.''
  \item \textbf{Khi bị hỏi ``có chắc luật đúng không''}: ``Luật quyết trong reducer
        thuần và có 23 test tự động trong \texttt{reducer.test.ts} kiểm chứng
        lifecycle, reload, Big Shot, thua theo vị trí và điều kiện thắng.''
\end{enumerate}

\end{document}
